\section{Experimental setup}
\label{sec:experiments}

\subsection{Corpora}
\label{sec:experiments:data}

Training draws from seven corpora, summarized in \Cref{tab:corpora}: \textbf{\texttt{data\_v5}} (evidence
-- SNLI \citep{bowman2015snli}, MNLI \citep{williams2018mnli}, ANLI \citep{nie2020anli}, BoolQ
\citep{clark2019boolq} -- plus wide-vocabulary classification including CLINC-150 and Banking77/HWU64
\citep{larson2019clinc,casanueva2020banking77}); \textbf{\texttt{data\_kb}} (closed-book multiple-choice
knowledge distilled from nine standard sources, TruthfulQA-MC1 \citep{lin2022truthfulqa} held out, leakage
audited against MMLU-Pro \citep{wang2024mmlupro,hendrycks2021mmlu}); \textbf{\texttt{data\_wf}} and
\textbf{\texttt{data\_wf\_hf}} (our rubric-conditioned workflow generator plus three externally sourced
typed-decision sets); \textbf{\texttt{data\_wf\_long}} (a longer-state, facts-first variant of the
workflow generator, \Cref{sec:method:mixture}); \textbf{\texttt{data\_wh}} (DecisionMix v2, a second
independent rubric-conditioned generator with compositional depth levels 1--7); and
\textbf{\texttt{data\_u}} (a soft-target uncertainty corpus from UNLI \citep{chen2020unli}, AmbiEnt, and
synthetic closed-form generators, with ChaosNLI-ambiguity \citep{nie2020chaosnli} evaluation-only). Every
corpus holds out whole families, styles, grammars, or parameter ranges as detailed in \Cref{tab:corpora};
full per-family breakdowns, licenses, and held-out split sizes are in \Cref{app:corpora}. None of the seven
corpora contains any of the 231 public JevBench ids.

\begin{table}[t]
  \centering
  \caption{Training corpora. ``Held out'' counts whole families, domains, styles, or parameter ranges
  excluded from training entirely (not a random validation split).}
  \label{tab:corpora}
  \small
  \begin{tabular}{lrrl}
    \toprule
    Corpus & Train rows & Held out & Role \\
    \midrule
    \texttt{data\_v5} & 633{,}213 & label wordings & Evidence (NLI/BoolQ) + wide-vocabulary classification \\
    \texttt{data\_kb} & 147{,}447 & TruthfulQA-MC1, GPQA & Closed-book multiple-choice knowledge \\
    \texttt{data\_wf} & 72{,}006 & 3 of 15 families & Rubric-conditioned workflow (Noul/Choice/Score) \\
    \texttt{data\_wf\_hf} & 86{,}600 & 3 external test/demo splits & External workflow sets (forms, rules, distillation) \\
    \texttt{data\_wf\_long} & 23{,}318 & -- & Long-state (multi-paragraph) workflow rows \\
    \texttt{data\_wh} & 60{,}605 & 1 domain/level, 3 domains, 2 styles, level 7 & DecisionMix v2 hard curriculum \\
    \texttt{data\_u} & 31{,}029 & 4 synthetic parameter ranges & Soft-target uncertainty corpus \\
    \bottomrule
  \end{tabular}
\end{table}

\subsection{External, evaluation-only sets}
\label{sec:experiments:external}

Six sets are never trained on and measure generalization beyond our own generators: \textbf{PagerDuty}
(alert-triage yes/no, floor .792), \textbf{jevlogs} (log-triage yes/no, floor .697), \textbf{Mind2Web}
(web-agent action choice, floor .427), \textbf{tree-choice} (taxonomy-routed choice up to $K{=}320$),
\textbf{typed-decisions} (400 cases spanning all three primitive types with full instructions, criteria,
and probabilistic gold), and \textbf{TruthfulQA-MC1}. We report accuracy against each set's own
majority-class floor because several are class-imbalanced enough that a constant prediction clears a
nontrivial accuracy; a model below its floor is worse than guessing the majority class.

\subsection{Held-out design}
\label{sec:experiments:heldout}

Generalization is measured along four largely independent axes, each with its own held-out split rather
than a random validation set: \emph{whole families} (an entire workflow family excluded from training),
\emph{rubric styles} (a writing style never seen in training), \emph{rule grammars} (one domain/level
combination's surface grammar withheld from DecisionMix v2), and \emph{compositional depth} (level 7:
temporal, unit-conversion, expected-value, or trade-off reasoning the generator never produces at any
trained level). We additionally report \textbf{rubric-flip accuracy} -- for a state and candidate set
rendered under two rubrics with different gold answers, the fraction of pairs where the model gets
\emph{both} answers correct, only possible if the rubric conditions the answer -- and
$\dr = \mathrm{Acc}(\state, r, \cands) - \mathrm{Acc}(\state, r_{\mathrm{shuf}}, \cands)$, comparing
accuracy under the correct rubric against a swapped one, state and candidates fixed.

\subsection{JevBench protocol and disclosures}
\label{sec:experiments:jevbench}

JevBench (\texttt{fstandhartinger/jevbench} v1.2.1) is a third-party benchmark for typed decision models,
with items typed \texttt{noul} (yes/no), \texttt{choice} ($K=2$--6), and \texttt{score} (ordinal), each
supplying a state, an instruction/rubric, and an exact label set, split into standard, easy, and hard
tiers. We report results with the following disclosures made explicit, because each materially affects
how the numbers should be read.

\textbf{Public subset, not a ranked entry.} The harness's public release exposes 231 of its items (72
standard, 48 easy, 111 hard); a further 146 judge-scored items are not public. The harness's own
leaderboard ranks only submissions covering at least 95\% of items including the private judge set. Every
number we report is computed on the 231 public ids only and is \emph{not} a ranked leaderboard entry;
comparisons to leaderboard entries in \Cref{sec:results} are drawn from the same public per-item file and
are only valid for that reason.

\textbf{Effective sample size.} With 72 standard items, the majority-class floor is .311 and the standard
error of a proportion at that floor is $\approx .058$ ($n{=}72$); design effects from paraphrase pairs in
the standard split further reduce the effective sample to $n_{\mathrm{eff}} = 36$. We treat any
standard-tier difference between our own checkpoints under roughly 2 points, and any hard-tier difference
under roughly 6 points, as noise rather than a ranking, and we say so explicitly wherever such a
comparison appears.

\textbf{Probabilities are conditioned on non-$\nullsym$.} The harness's scoring expects a distribution
over the labeled candidates only; we report $P(\cand{j} \mid \state, \query, \cands, \neg\nullsym)$, i.e.\
$\pnull$ is computed and then dropped and the remaining mass renormalized, rather than folding $\pnull$
into the reported distribution.

\textbf{Other disclosures}, held constant across every model we report: option order follows the
harness's own label order (a reversed-order control on one checkpoint swung standard accuracy by $\pm 4$
points, larger than the differences between our own model variants); latency is measured in-process on
one H100, one decision at a time, with model loading excluded and cold candidate-label embedding included
where applicable; and no hosted-API cost applies to any of our numbers (we are not a hosted service on
this benchmark). Every run's harness commit, checkpoint hash, and dependency lock file are recorded in
the run manifest.

\subsection{Metrics}
\label{sec:experiments:metrics}

We report \textbf{accuracy} (argmax against gold), \textbf{Brier score}, \textbf{expected calibration
error (ECE)}, and \textbf{negative log-likelihood (NLL)}, against hard or soft gold as the evaluation set
provides. Where a label space is only partially covered by the candidate set at test time, we separately
report \emph{among-$K$} accuracy (argmax restricted to the true candidate set, ignoring $\pnull$) so null
and discrimination failures are not conflated. $\dq$ is the gap between a model's score on a question and
its score with the question replaced by a placeholder or shuffled one (candidate set fixed), isolating
question-conditioned signal from candidate-set priors; we report both a choices-only and a stricter
shuffled-question variant, the latter primary after the former showed residual candidate-prior leakage
(\Cref{sec:method:blind-closure}). $\dr$ and rubric-flip accuracy are defined in
\Cref{sec:experiments:heldout}. Rubric-conditioned score and typed external sets carry soft, non-one-hot
gold, so NLL and Brier score are primary there and accuracy is a secondary, coarser signal.

\subsection{Hardware and latency measurement}
\label{sec:experiments:hardware}

Training runs on rented NVIDIA H100 GPUs (SXM and NVL variants), bf16, AdamW, effective batch size 64,
8{,}000--12{,}000 steps depending on the run. Latency is measured in-process, single stream, one decision
at a time unless a batch size is stated, with the state's key-value cache warm and model loading excluded;
we report single-decision latency and marginal per-query latency at a stated $M$, both as a function of
$K$. Numbers compared across backbone sizes in one table are measured on one pod with one PyTorch/CUDA
build (``apples-to-apples''); numbers from different pods or builds are marked as such and not compared
directly.
